% !TEX root = ../main.tex
\section{Contextual Truth and Illicit Universalization}
\label{sec:contextual-truth}

Before defining the grounding operator I need one more piece of apparatus, and
it is the piece that practitioners in regulated domains will recognise first. It
concerns truth rather than assertion, and it is what forces the context index
that Section~\ref{sec:grounding} builds into $\Grd_M$.

Many propositions that look simple in natural language are incomplete until
their governing context is fixed. In regulated, clinical, legal, and financial
settings, truth is routinely relative to jurisdiction, time, purpose, and the
facts of a particular matter. Accordingly the truth condition for a practical
proposition should be indexed:
\[
  \True(\phi \mid j, \tau, u, f),
\]
where $j$ is jurisdiction, $\tau$ time, $u$ intended use, and $f$ the relevant
facts. I will write $c = (j, \tau, u, f)$ for the resulting practical context and
use $\True(\phi \mid c)$ where the components need not be displayed.

A proposition concerning legal permissibility, clinical appropriateness,
accounting treatment, or standard of care can change truth value when any
component changes. A claim about an approved indication may be true in one
jurisdiction and false in another. A claim true under a prior regulation may be
false after amendment. A claim appropriate for one patient population may be
inappropriate for another. Nothing here is exotic; it is the ordinary condition
of consequential work.

The failure mode this creates is not primarily factual inaccuracy. It is
\emph{unresolved reference}. A system may produce a sentence with apparent
propositional form while leaving its decisive parameters unspecified. Ask whether
something is approved, or compliant, or permitted, and the question as posed
typically leaves open the jurisdiction, the effective date, the population, the
governing rule, the intended action, and the risk tolerance. Let
\[
  \Resolved(\phi, c)
\]
mean that the material context required to interpret $\phi$ has been fixed. Then
a consequential claim should not be treated as fully evaluable unless its context
is resolved:
\[
  \Evaluable(\phi, c) \Rightarrow \Resolved(\phi, c).
\]

Failure to resolve creates what I will call \emph{illicit universalization}: the
error of treating a proposition conditionally true in one context as true
simpliciter. Formally, the entailment that fails is
\[
  \True(\phi \mid j, \tau, u, f)
  \nentails
  \forall j' \forall \tau' \forall u' \forall f'\,
  \True(\phi \mid j', \tau', u', f').
\]
In prose: the truth of a proposition under one specified context does not entail
its truth across all contexts. Note the direction of the claim. It is not that
differing contexts \emph{guarantee} differing truth values---a proposition may
well be true under both---but that context-relative truth licenses no
generalisation. The weaker claim is the one that does the work, and the stronger
one is false.

This bears on the argument of Sections~\ref{sec:output-operator}
and~\ref{sec:hyperintensionality} in a way worth making explicit. Illicit
universalization is what one should expect from a hyperintensional operator
computed over formulations. A formulation carries no index; the practical context
that would make its content evaluable is not part of the string, and there is no
mechanism whose function is to notice that it is missing. A system that
represented propositions would face the question of which proposition an
underspecified query expresses. A system that operates on formulations does not
face it, and so produces a confident answer to a question that has not yet been
asked.
